\section{\change{Background Details}}


\subsection{\change{background about current memory systems}}
\label{background_detail}
\change{We describe both the mainstream memory architectures and the \textbf{LightMem} pipeline in terms of two major stages. 
The first is the memory bank construction stage, which can be further decomposed into the three sub-stages (I), (II), and (III) described in the Section~\ref{main_text_preliminary}. 
The second major stage concerns the usage of the memory system, which consists of retrieval and question answering (QA). }

\paragraph{\change{Memory Bank Construction}}
\change{As shown in Table~\ref{tab:detail_process}, we detail the workflows of the three sub-stages (I), (II), and (III) for naive RAG, prevailing memory systems, and our LightMem. 
It can be observed that baseline memory systems typically perform their update stage during user–model interaction, which introduces substantial test-time latency. 
In contrast, LightMem decouples this update process from online interaction, thereby significantly reducing test-time latency.
All models involved in these processes are listed in Table~\ref{tab:function_model_map}. 
As shown, LightMem introduces only one additional model, LLMlingua-2,beyond those used by baseline methods. 
This model follows a lightweight BERT architecture and requires less than 2GB of GPU memory during inference, rendering its overhead negligible. 
Moreover, for fairness, the latency introduced by this component is fully accounted for in our reported Runtime metric.}

\begin{table*}[ht]
\centering
\caption{\change{The mainstream memory architectures and the LightMem pipeline of memory bank construction stage. 
Black-font processes denote those executed during online test-time interactions, whereas \textcolor{red}{red-font} processes denote those executed offline.
}}
\vspace{0.5ex}
\renewcommand{\arraystretch}{1.1}
\setlength{\tabcolsep}{3pt}
\small
\scalebox{0.95}{
\begin{threeparttable}
\begin{tabular}{@{}lccc@{}}
\toprule
\textbf{Method} & \textbf{(I) Segment} & \textbf{(II) Summary/Extrct} & \textbf{(III) Update}\\ 
\midrule
NaiveRAG
& \makecell{$ \text{Raw dialog} \to f_{\text{seg}}()$ \\ 
$\to \{ \text{seg}_i \} $}

& $ \to f_{\text{index}}() \to \{\text{emb}_i\} $
& \textbackslash
\\
\midrule
\makecell{Other \\ Memory \\ Systems}
& \makecell{$ \text{Raw dialog} \to f_{\text{seg}}()$ \\ 
$\to \{ \text{seg}_i \} $}
& \makecell{$ \to f_{\text{sum/extract}}() \to \{\text{memory entry}_i\} $ \\
$ \to f_{\text{index}}() \to \{\text{emb}_i\} $
}
& \makecell{$ \to f_{\text{retrieve}}() \to \{\text{related entry}_i\} $ \\
$ \to f_{\text{update}}() $ \\ 
$  \to \{\text{add, delete, update, merge...}\} $
}
\\
\midrule
LightMem 
& \makecell{$\text{Raw dialog} \to f_{\text{seg}}()$ \\ 
$ \to \{\text{seg}_i\}$ \\
$\to f_{\text{pre\_compress}}() $ \\
$\to \{\text{comp\_seg}_i\}$ \\
$\to \text{sensory buffer full} \to$ \\ 
$f_{\text{topic}}() \to $ \\
$\{\text{topic-wise comp\_seg}_i\}$
}

& \makecell{$\to \text{STM buffer full} $ 
$\to f_{\text{sum/extract}}()$ \\
$\to \{\text{topic}_i, \{\text{memory entry}_j\}\}$ \\
$ \to f_{\text{index}}() \to \{\text{topic}_i, \{\text{emb}_j\}\} $
}
& \makecell{\color{red}$\text{Offline update trigger} $ \\
\color{red}$\{\text{every entry}_i\} \to f_{\text{retrieve}}()$ \\
\color{red}$\to \{\text{related entry}_j\} \to \{\text{update queue}\}$ \\
\color{red}All update queues established \\
\color{red}$ \to \text{parallel } f_{\text{update}}() $ \\ 
\color{red}$ \to \{\text{add, delete, update, merge...}\} $
}

\\


\bottomrule
\end{tabular}
\end{threeparttable}
}
\label{tab:detail_process}
\end{table*}

\begin{table}[ht]
\centering
\renewcommand{\arraystretch}{1.2}
\begin{tabular}{l l l}
\hline
\textbf{Function} & \textbf{Model / Strategy} & \textbf{Implementation in This Paper} \\
\hline
$f_{\text{seg}}()$ & Segmentation strategy & Turn-level granularity input \\
$f_{\text{index}}()$ & Embedding model & all-MiniLM-L6-v2 \\
$f_{\text{sum/extract}}()$ & System backbone model & GPT-4o-mini; Qwen3-30B-A3B-Instruct-2507 \\
$f_{\text{retrieve}}()$ & Retrieval strategy & Cosine similarity vector retrieval \\
$f_{\text{update}}()$ & System backbone model & GPT-4o-mini; Qwen3-30B-A3B-Instruct-2507 \\
\color{red}$f_{\text{pre\_compress}}()$ & \color{red}Token compression model & \color{red}LLMlingua-2 \\
\color{red}$f_{\text{topic}}()$ & \color{red}Topic segmentation model & \color{red}LLMlingua-2 \\
$f_{\text{chat}}()$ & Chat model & GPT-4o-mini; Qwen3-30B-A3B-Instruct-2507 \\
\hline
\end{tabular}
\caption{\change{Mapping between functions, their roles, and the concrete models used in this paper. Black-font entries denote models shared by both LightMem and baseline methods, whereas \textcolor{red}{red-font} entries denote models unique to LightMem.}}
\label{tab:function_model_map}
\end{table}

\paragraph{\change{Retrieval and Usage}}
\change{After the memory bank construction stage, we obtain an up-to-date memory bank. 
When a new user query arrives, the memory system use $f_{\text{retrieve}}()$ to retrieve relevant entries from this repository, appends them to the query, and then prompts the chat model $f_{\text{chat}}()$ to produce a response.} 


% \subsection{Metric Details}

% \label{metric_detail}

\subsection{\change{Notation and Complexity Details}}
\label{notation_detail}

\begin{table}[h]
\centering
\caption{\change{Notation used in complexity analysis (\S Section~\ref{sec:complexity_analysis}).}}
\begin{tabular}{p{2cm} p{12cm}}
\toprule
\textbf{Symbol} & \textbf{Definition} \\
\midrule
$N$ & Total number of turns in a dialogue history. \\
\midrule
$T$ & Average number of tokens per turn. \\
\midrule
$r$ & Token compression rate (as defined in the main paper). After one compression step, only a fraction $r$ of tokens is retained. \\
\midrule
$x$ & Number of compression iterations. In LightMem, the \textit{pre-compress} module may be invoked multiple times for the same message to remove redundancy until the message is sufficiently compact. This occurs frequently in datasets such as \textbf{LongMemEval}. All time costs are included in runtime metrics. \\
\midrule
$th$ & Capacity of the Short-Term Memory (STM) buffer, as defined in the paper. \\
\midrule
\makecell{$L_{\text{sum-in}}$ / $L_{\text{sum-out}}$} & Number of tokens in the \textbf{input prompt template} and \textbf{output} of a single backbone LLM call for \textit{summarization}. These are similar across memory frameworks. \\
\midrule
$M_1$ / $M_2$ & Number of memory entries produced from a single summarization operation under Other Memory Systems ($M_1$) and LightMem ($M_2$). \\
\midrule
$L_{\text{up-in}}$ / $L_{\text{up-out}}$ & Number of tokens in the \textbf{input prompt template} and \textbf{output} of a single backbone LLM call for \textit{memory update}. Similar across frameworks. \\
\midrule
$R_1$ / $R_2$ & Proportion of summary entries that successfully retrieve at least one relevant memory entry (triggering an update) for Other Memory Systems ($R_1$) and LightMem ($R_2$). Some entries do not retrieve any relevant counterparts and thus do not trigger updates. \\
\bottomrule
\end{tabular}
\label{tab:notation}
\end{table}

